% Part: methods
% Chapter: proofs
% Section: introduction

\documentclass[../../../include/open-logic-section]{subfiles}

\begin{document}

\olfileid[tr]{mth}{prf}{int}

\olsection{Giriş}

Başlangıç mantığı derslerindeki deneyiminiz sayesinde, muhtemelen bir
!!a{derivation} sistemiyle---büyük olasılıkla bir doğal türetim ya da
Fitch tarzı !!{derivation} sistemiyle, belki de bir kanıt ağacı
sistemiyle---rahat çalışıyorsunuzdur. Bu sistemlerde kanıtlar yaparak ya
bir !!a{formula} kanıtladığınızı ya da verilen bir argümanın geçerli
olduğunu gösterdiğinizi muhtemelen hatırlıyorsunuzdur. Bunu yapmak için,
istenen sonuca ulaşıncaya kadar sistemin kurallarını uyguladınız. Mantık
\emph{hakkında} akıl yürütürken de birtakım şeyleri kanıtlarız; ancak
çoğu durumda bir !!a{derivation} sistemi kullanmayız. Aslında ele
aldığımız kanıtların çoğu, bütünüyle birinci dereceden mantığın dilinde
değil, Türkçe olarak (belki araya biraz simgesel dil katılarak)
yapılır. Böyle kanıtlar kurarken başlangıçta ne yapacağınızı
bilemeyebilirsiniz---bir !!a{derivation} sistemi olmadan bir şeyi nasıl
kanıtlarım? Nereden başlarım? Kanıtımın doğru olduğunu nasıl anlarım?

Bir kanıta girişmeden önce kanıtın ne olduğunu ve nasıl kurulacağını
bilmek önemlidir. Adından da anlaşılacağı üzere bir \emph{kanıt}, bir
şeyin doğru olduğunu göstermeye yarar. Bunu bir diyalog biçiminde
düşünebilirsiniz---birisi size bir şeyin doğru olup olmadığını, örneğin
iki dışındaki her asal sayının tek olup olmadığını sorar. Yalnızca
``evet'' demek yetmez; \emph{nedenini} bilmek isteyebilir. Bu durumda ona
bir kanıt verirsiniz.

Gündelik konuşmada bir yanıta kabaca işaret etmek ya da eksik bir yanıt
vermek yeterli olabilir. Oysa mantıkta ve matematikte titiz bir kanıt
isteriz---bir şeyin \emph{hiçbir} kuşkuya yer bırakmayacak biçimde doğru
olduğunu göstermek isteriz. Bu, kanıtımızdaki her adımın
gerekçelendirilmesi ve gerekçenin sağlam olması gerektiği anlamına gelir
(yani kullandığınız varsayımın gerçekten kanıtladığınız teoremin
ifadesinde varsayılmış olması, uyguladığınız tanımların doğru
uygulanması, başvurduğunuz gerekçelerin doğru çıkarımlar olması vb.).

Genellikle bir ifadeyi kanıtlarız. Kanıtladığımız ifadelere çeşitli
adlar veririz: önermeler, teoremler, lemmalar ya da sonuçlar. Önerme,
kanıtlanmaya değer temel bir ifadedir: kayda geçirilecek kadar önemlidir
ama özellikle derin olmayabilir veya sık kullanılmayabilir. Teorem,
önemli ve dikkate değer bir önermedir. Kanıtı çoğu zaman birkaç adıma
ayrılır; bazen ilk kanıtlayan kişinin (örneğin Cantor Teoremi,
L\"owenheim--Skolem teoremi), bazen de ilgili olduğu olgunun (örneğin
tamlık teoremi) adını taşır. Lemma, daha önemli bir sonucun kanıtında
kullanılan bir önerme ya da teoremdir. Kafa karıştırıcı biçimde,
lemmalar bazen başlı başına önemli sonuçlardır ve onları ortaya koyan
kişinin adını da taşırlar (örneğin Zorn Lemması). Sonuç ise başka bir
sonuçtan kolayca çıkan bir ifadedir.

Kanıtlanacak ifade çoğu kez, ne tür şeyler hakkında bir şey
kanıtladığımızı açıklığa kavuşturan varsayımlar içerir. İfade ``$!A$,
$!B \lif !C$ biçiminde bir !!a{formula} olsun'' ya da ``$\Gamma
\Proves !A$ olduğunu varsayalım'' veya buna benzer bir sözle
başlayabilir. Bunlar önermenin, teoremin ya da lemmanın
\emph{hipotezleridir}; kanıtınızda bunların doğru olduğunu
varsayabilirsiniz. Bunlar kanıtladığımız şeyi sınırlar ve ayrıca
hakkında konuştuğumuz nesneler için birtakım adlar ortaya koyar.
Örneğin önermeniz ``$!A$, $!B \lif !C$ biçiminde bir !!a{formula}
olsun'' diye başlıyorsa, yalnızca belirli türdeki bütün formüller
(yani koşullu önermeler) hakkında bir şey kanıtlıyorsunuzdur ve
$!B \lif !C$'nin, kanıtınızın ele alacağı keyfî bir koşullu önerme
olduğu anlaşılır.

\end{document}
